% !TEX root = main.tex
\documentclass[main.tex]{subfiles}
\begin{document}

\section{Site optimization}
\label{sec:site-min}

So far the sites $\qq$ have been fixed inputs to which the partition adapted.
However, sites are parametric handles for the power diagram and do not carry direct geometric meaning, whereas the cell centroids encode the first-order position of each cell.
It is therefore natural to express design objectives in terms of centroids rather than sites.
As shown in \S\ref{sec:remarks}, the site gradient of the transport cost vanishes when $\qq_i = \tcentroid_i$, so optimizing over the sites at fixed targets places each site at its prescribed target; the sites thus coincide with the centroid coordinates at optimality.

We introduce a distortion potential $\distortion \colon \R^{n \times d} \to \R$ measuring how far the target centroids deviate from a desired configuration.
Combining it with the transport cost, we seek the sites, target centroids, and partition that jointly solve
\begin{subequations}
  \label{eq:site-full-problem}
  \begin{alignat}{2}
    & \min_{\qq,\, \tcentroid,\, \cell} \quad && \tau\,\distortion(\tcentroid) + \sum_{i} \int_{\cell_i} \norm{\xx - \qq_i}^2 \, d\xx \label{eq:site-full-problem-obj} \\
    & \text{s.t.} \quad && \cellvol_i = \pvol_i, \ \ \int_{\cell_i} (\xx - \tcentroid_i) \, d\xx = 0 \ \ \forall i , \label{eq:site-full-problem-cons}
  \end{alignat}
\end{subequations}
where $\tau > 0$ controls the strength of the distortion relative to the transport cost.
Since $\distortion$ depends only on $\tcentroid$ and the site gradient of the transport cost vanishes at $\qq_i = \tcentroid_i$ (\S\ref{sec:remarks}), the optimal sites coincide with the targets.
Setting $\tcentroid = \qq$ and invoking the centroidal joint dual $\barjointdual$ from \S\ref{sec:remarks}, problem~\eqref{eq:site-full-problem} reduces to the saddle-point problem
\begin{equation}
  \min_{\qq} \max_{\shift,\, \weight} \
    \tau\,\distortion(\qq) + \barjointdual(\shift, \weight) ,
  \label{eq:site-saddle}
\end{equation}
for given target volumes $\pvol_i$, coupling the outer minimization in $\qq$ with the concave inner maximization of \S\ref{sec:joint-constraints}.
Since $\tcentroid_i = \qq_i$ and $\cellvol_i = \pvol_i$ at the inner optimum, the conservation identity~\eqref{eq:partition-feasibility-centroid} imposes a feasibility constraint on the sites,
\begin{equation}
  \sum_i \pvol_i \, \qq_i = \domainVol \, \centroid_\domain .
  \label{eq:site-feasibility}
\end{equation}
Our final formulation is thus
\begin{equation}
  \min_{\qq} \max_{\shift,\, \weight} \
    \tau\,\distortion(\qq) + \barjointdual(\shift, \weight)
  \quad \text{s.t.} \quad
    \sum_i \pvol_i \, \qq_i = \domainVol \, \centroid_\domain ,
  \label{eq:site-saddle-feasible}
\end{equation}
a saddle-point problem that can be solved by alternating between the maximization over the multipliers $(\shift, \weight)$ for fixed sites and the constrained minimization over the sites $\qq$ for fixed multipliers.

\paragraph{Gradient.}
For the site minimization at fixed multipliers, the centroidal dual $\barjointdual$ depends on $\qq$ through the shifted sites $\dual_i = \qq_i + \shift_i$.
Differentiating the centroidal Lagrangian at the inner-optimal partition gives the site gradient
\begin{equation}
  \partial_{\qq_i} \barjointdual
    = -2 \int_{\cell_i} (\xx - \qq_i) \, d\xx + 2\,\pvol_i\,\shift_i
    = 2\,\pvol_i\,\shift_i ,
  \label{eq:site-centroidal-gradient}
\end{equation}
where the transport term vanishes by centroid feasibility ($\centroid_i = \qq_i$).
Including the distortion and enforcing the feasibility constraint~\eqref{eq:site-feasibility} via a multiplier $\yy \!\in\! \R^d$, stationarity of the site minimization reads, per site,
\begin{equation}
  \tau\,\partial_{\qq_i} \distortion(\qq) + 2\,\pvol_i\,\shift_i - \pvol_i\,\yy = 0 .
  \label{eq:site-stationarity}
\end{equation}

\paragraph{Quadratic distortion.}
When $\distortion$ is a convex quadratic of the form
\begin{equation}
  \distortion(\qq)
    = \tfrac{1}{2}\,\operatorname{tr}\!\bigl(\qq\tr \Amat\, \qq\bigr)
    + \operatorname{tr}\!\bigl(\qq\tr \bmat\bigr) ,
  \label{eq:site-quadratic}
\end{equation}
with $\Amat \!\in\! \R^{n \times n}$ symmetric positive semi-definite and $\bmat \!\in\! \R^{n \times d}$, the per-site gradient is $\partial_{\qq_i} \distortion = \sum_j \Amat_{ij}\, \qq_j + \bmat_i$, and \eqref{eq:site-stationarity} becomes
\begin{equation}
  \tau \Bigl( \sum_j \Amat_{ij}\, \qq_j + \bmat_i \Bigr) + 2\,\pvol_i\,\shift_i - \pvol_i\,\yy = 0 .
  \label{eq:site-stationarity-quadratic}
\end{equation}
For fixed multipliers, this is a linear system in the sites $\qq$, coupling each site to its neighbors through $\Amat$.

\end{document}
